% Introduction to the EU ... Week 9 Housing Crisis Workshop
% Compile: pdflatex eu-housing-crisis-workshop.tex

\documentclass[10pt,a4paper]{article}
\usepackage[utf8]{inputenc}
\usepackage[T1]{fontenc}
\usepackage[margin=1.2cm]{geometry}
\usepackage{enumitem}
\setlist{nosep,leftmargin=*}
\usepackage{xcolor}
\usepackage{fancyhdr}
\usepackage{microtype}
\usepackage{hyperref}

\definecolor{eublue}{RGB}{0,51,153}
\pagestyle{fancy}
\fancyhf{}
\fancyfoot[L]{\small Intro EU | Week 9}
\fancyfoot[R]{\small\thepage}
\raggedright

\hypersetup{
  colorlinks=true,
  urlcolor=eublue,
}

\begin{document}
\begin{center}
{\large\bfseries\color{eublue} Introduction to the EU ... Week 9}\\[0.2em]
\footnotesize Mar 23. Housing crisis workshop (European Parliament ... group debate)
\end{center}\vspace{0.5em}

\section*{\color{eublue}I. Workshop Goal --- Housing Crisis Overview}
\begin{enumerate}[label=\Alph*.]
\item \textbf{Work focus}
  \begin{enumerate}[label=\arabic*.]
  \item Housing crisis in Europe
  \item Latest initiatives (very recent)
  \item Compare positions / roots of arguments
  \end{enumerate}
\item \textbf{Main outputs (group)}
  \begin{enumerate}[label=\arabic*.]
  \item Understand opinions expressed
  \item Identify lexicon / words used
  \item Identify different positions
  \item Note similarities vs differences + why
  \end{enumerate}
\item \textbf{Evidence sources (given)}
  \begin{enumerate}[label=\arabic*.]
  \item Europarl website search: housing
  \item Videos + recorded materials (as indicated)
  \item Statements / debates on housing in parliament
  \item EU group roles + what groups represent
  \end{enumerate}
\end{enumerate}

\section*{\color{eublue}Session readings (week-9 PDFs)}
\begin{itemize}
\item \texttt{7a Parlamentary Europe.pdf}
\item \texttt{7b Olsen European Parliament.pdf}
\item \texttt{European Parliament.pdf}
\end{itemize}

\section*{\color{eublue}Materials link (for reference)}
\begin{itemize}
\item Europarl housing search:
  \href{https://www.europarl.europa.eu/portal/en/search?planet=_all&searchQuery=housing}{https://www.europarl.europa.eu/portal/en/search?planet=_all\&searchQuery=housing}
\end{itemize}

\section*{\color{eublue}Housing focus checklist (for group work)}
\begin{itemize}
\item Use recent housing material from Europarl (statements, debates, videos, initiatives).
\item For each group, note: \textbf{housing lexicon}, \textbf{position}, and \textbf{main policy tool}.
\item Compare groups: \textbf{similarities}, \textbf{differences}, and short reasons for each difference.
\item Keep wording close to what they actually say in the source.
\end{itemize}


\section*{\color{eublue}II. Activity Instructions --- How to Work}
\begin{enumerate}[label=\Alph*.]
\item \textbf{Steps (as described)}
  \begin{enumerate}[label=\arabic*.]
  \item Go to Europarl website --- use search (housing)
  \item Check latest initiatives responding to crisis
  \item Watch recorded/video examples (if provided)
  \item Identify positions expressed by groups
  \item Prepare in groups for short presentation
  \end{enumerate}
\item \textbf{Planning / timing}
  \begin{enumerate}[label=\arabic*.]
  \item Group work first
  \item Prepare presentation for class
  \item 45 minutes group work (as mentioned)
  \end{enumerate}
\end{enumerate}

\section*{\color{eublue}III. Institutions and Actors --- Parliamentary Debate Lens}
\begin{enumerate}[label=\Alph*.]
\item \textbf{European Parliament (frame)}
  \begin{enumerate}[label=\arabic*.]
  \item Use parliament vocabulary (as indicated)
  \item Groups in parliament $\rightarrow$ represent positions
  \item Role-play running order: ``President'' manages the debate flow (mini-plenary)
  \end{enumerate}
\item \textbf{Mentioned admin/political dimension (from transcript)}
  \begin{enumerate}[label=\arabic*.]
  \item Ministers / ministries referenced
  \item Member states / partners referenced
  \item Procedure logic: if the proposal can proceed / ``pass'' in the parliamentary flow $\rightarrow$ continue; otherwise stop / skip reading
  \end{enumerate}
\end{enumerate}

\section*{\color{eublue}IV. Case Example (Historical Policy) --- Vienna Housing}
\begin{itemize}
\item \textbf{Prevention logic}: avoid secondary/undesired effects.
\item \textbf{Public housing requirement scale} (as in class example): approx. 1940s $\sim$50\%; center areas $\sim$70\% to $\sim$50\%.
\item \textbf{Effects}: price + family access (what changes for who?)
\item \textbf{Why it matters for today}: real-world effect when searching for a house
\item \textbf{Complexity / difficult demand context}: as described in class example
\end{itemize}

\section*{\color{eublue}V. Presentation and Debate Mechanics}
\begin{enumerate}[label=\Alph*.]
\item \textbf{In-class flow}
  \begin{enumerate}[label=\arabic*.]
  \item Groups present prepared statements/activities
  \item Listen to other groups' positions
  \item Debate based on statements
  \end{enumerate}
\item \textbf{Debate constraints / rules (as described)}
  \begin{enumerate}[label=\arabic*.]
  \item Use of selected statements
  \item Compare lexicon + roots
  \item Similarities + differences + reasons
  \item President flow: if a proposal can pass, continue; if not, stop/skip
  \end{enumerate}
\end{enumerate}

\section*{\color{eublue}Political groups to compare (10th EP)}
\begin{itemize}
\item EPP
\item S\&D
\item Renew
\item Greens/EFA
\item The Left
\item ECR
\item PfE
\item ESN
\item Non-Inscrits
\end{itemize}

\section*{\color{eublue}Required source log (one Europarl item per group)}
\begin{itemize}
\item For each group you discuss, you must paste the exact Europarl link you used (debate/initiative + date + speaker if you can).
\end{itemize}
\begin{center}
\begin{tabular}{p{2.6cm} p{10.5cm}}
\textbf{Group} & \textbf{Europarl link (paste URL + item title)} \\
\hline
EPP & \texttt{https://www.europarl.europa.eu/news/en/press-room/20260205IPR33603/meps-adopt-proposals-to-tackle-europe-s-housing-crisis} (Borja Gim\'enez Larraz, 9 Feb 2026) \\
S\&D & \texttt{https://www.europarl.europa.eu/news/en/press-room/20260205IPR33603/meps-adopt-proposals-to-tackle-europe-s-housing-crisis} (Aodh\'an \'O R\'iord\'ain, S\&D housing negotiator) \\
Renew & \texttt{https://www.europarl.europa.eu/news/en/press-room/20260205IPR33603/meps-adopt-proposals-to-tackle-europe-s-housing-crisis} (Ciar\'an Mullooly, Renew, Vice-Chair of HOUS) \\
Greens/EFA & \texttt{https://www.europarl.europa.eu/news/en/press-room/20260205IPR33603/meps-adopt-proposals-to-tackle-europe-s-housing-crisis} (Benedetta Scuderi, Greens) \\
The Left & \texttt{https://www.europarl.europa.eu/news/en/press-room/20260205IPR33603/meps-adopt-proposals-to-tackle-europe-s-housing-crisis} (Leila Chaibi / The Left housing justice framing) \\
ECR & \texttt{https://www.europarl.europa.eu/news/en/press-room/20260205IPR33603/meps-adopt-proposals-to-tackle-europe-s-housing-crisis} (ECR national competence framing) \\
PfE & \texttt{https://www.europarl.europa.eu/news/en/press-room/20260205IPR33603/meps-adopt-proposals-to-tackle-europe-s-housing-crisis} (PfE national sovereignty framing) \\
ESN & \texttt{https://www.europarl.europa.eu/news/en/press-room/20260205IPR33603/meps-adopt-proposals-to-tackle-europe-s-housing-crisis} (ESN sovereignty/bureaucracy framing) \\
Non-Inscrits & \texttt{PV-10-2026-03-10 (paste HOUS minutes link) + speaker name + item/quote title} \\
\end{tabular}
\end{center}

\section*{\color{eublue}VI. Political group stance snapshot (reference)}
\begin{itemize}
\item \textbf{EPP}: supply, investment, simplification; market-led tools and permit/tax easing; EU as facilitator (not strong market regulator).
\item \textbf{S\&D}: affordability, social emergency, tenants' rights; public/affordable housing + tenant protections; strong pro-EU intervention.
\item \textbf{Renew}: investment gap, affordability, first-time buyers; tax/investment instruments and market efficiency; pro-EU coordination.
\item \textbf{Greens/EFA}: right to housing, anti-speculation, public housing; Housing First and stronger anti-speculation rules; strong pro-EU intervention.
\item \textbf{The Left}: financialisation, social right, rent caps; public/cooperative housing and stronger tenant safeguards; pro-EU but critical of EU fiscal constraints.
\item \textbf{ECR}: subsidiarity, national competence, simplification; deregulation and national/local decisions; skeptical of binding EU housing regulation.
\item \textbf{PfE}: national sovereignty, family burden, energy costs; national control and lower household burden framing; very skeptical of EU housing mandates.
\item \textbf{ESN}: sovereignty, anti-bureaucracy, national priorities; redirect funding to national/regional priorities; very skeptical/hostile to EU-level housing regulation.
\item \textbf{Non-Inscrits}: no unified line; stance depends on each individual MEP.
\end{itemize}

\section*{\color{eublue}Political groups worksheet (simple)}
\begin{itemize}
\item Political groups to compare: \texttt{EPP}, \texttt{S\&D}, \texttt{Renew}, \texttt{Greens/EFA}, \texttt{The Left}, \texttt{ECR}, \texttt{PfE}, \texttt{ESN}, \texttt{Non-Inscrits}.
\item For each group you discuss (use one recent Europarl item + their wording):
  \begin{itemize}
  \item Housing lexicon: 3--5 key terms + 1 sentence defining the problem as they frame it.
  \item Policy tools: 1--2 emphasized mechanisms.
  \item EU role stance: pro-EU / cautious / skeptical (write one phrase they use).
  \item Debate logic: what should pass vs what should be blocked/skipped (based on their wording).
  \end{itemize}
\item Finish with comparison: \textbf{3 similarities} + \textbf{3 differences}; for each difference add \textbf{1 reason} tied to their wording/tools.
\end{itemize}

\section*{\color{eublue}VII. Extra institutional recap (from transcript) --- Democratic deficit}
\begin{itemize}
\item Democratic deficit (plain words): the EU can make big decisions, but citizens don't always feel they can directly influence those decisions.
\item Early integration: decisions were shaped largely via national contexts and big institutions (especially the Council); Commission/Council had strong role, while EP was supposed to represent citizens but had limited power.
\item Why it matters for today's debate lens: connects ``institutions + who controls law-making'' to workshop tasks (parliament vocabulary, procedures, groups).
\end{itemize}

\end{document}

